\chapter{Hyrje në simulimet fizike dhe programimin shkencor}
\begin{qellime}
Studenti dallon sistemin real, modelin matematik, algoritmin dhe kodin; formulon hyrjet, daljet dhe supozimet; dallon simulimin determinist nga ai stokastik; organizon një eksperiment të riprodhueshëm në Python ose MATLAB/Octave.
\end{qellime}

\section{Nga modeli te eksperimenti numerik}
Një simulim fizik është një eksperiment i kryer mbi një model. Ai nuk fillon me kodin, por me një pyetje: cilën madhësi kërkojmë, në çfarë kushtesh dhe me çfarë saktësie? Sistemi real përmban më shumë hollësi sesa mund të përfaqësohen. Modelimi zgjedh mekanizmat dominues, ndërsa simulimi zbaton pasojat e kësaj zgjedhjeje.

Një problem dinamik mund të shkruhet në formën
\begin{equation}
 \frac{\dd \vect y}{\dd t}=\vect f(t,\vect y;\vect\theta),\qquad \vect y(t_0)=\vect y_0,
\end{equation}
ku $\vect y$ është gjendja dhe $\vect\theta$ përmbledh parametrat. Për një grimcë në një përmasë, $\vect y=(x,v)^T$. Në një model stokastik, evolucioni varet edhe nga një ndryshore rasti $\xi_n$,
\begin{equation}
 \vect Y_{n+1}=\vect G(t_n,\vect Y_n;\vect\theta,\xi_n).
\end{equation}
Një ekzekutim determinist jep të njëjtën trajektore për të njëjtat hyrje; një ekzekutim stokastik jep një realizim. Përfundimi stokastik merret nga bashkësia e realizimeve.

\section{Hierarkia e modeleve}
Modelet ndërtohen në nivele. Modeli minimal izolon mekanizmin kryesor dhe lehtëson kontrollin analitik. Modeli i zgjeruar shton fërkim, forcë të jashtme, zhurmë ose gjeometri reale. Modeli më kompleks nuk është domosdoshmërisht më i mirë: ai kërkon më shumë parametra, të dhëna dhe prova vleftësimi.

\begin{shembull}[Rënia e lirë]
Pa rezistencë ajri, $y(t)=y_0+v_0t-gt^2/2$. Ky rast ka zgjidhje analitike dhe përdoret për verifikim. Me rezistencë lineare, $m\dot v=-mg-cv$; me rezistencë kuadratike, $m\dot v=-mg-cv|v|$. Zgjedhja ndërmjet tyre është vendim fizik, jo programues.
\end{shembull}

\section{Diskretizimi}
Kompjuteri ruan një numër të fundmë vlerash. Koha zëvendësohet nga rrjeti
\begin{equation}
 t_n=t_0+n\Delta t,\qquad n=0,1,\ldots,N.
\end{equation}
Hapi $\Delta t$ kontrollon rezolucionin dhe koston. Një hap më i vogël zakonisht ul gabimin e këputjes, por rrit numrin e veprimeve dhe mund të ekspozojë gabimin e rrumbullakimit. Përzgjedhja e hapit duhet të mbështetet në një studim konvergjence.

\section{Programimi shkencor}
Programimi shkencor e kthen arsyetimin në një procedurë që mund të përsëritet. Një notebook i mirë përmban: pyetjen fizike; supozimet; parametrat dhe njësitë; funksionet; eksperimentet; figurat; analizën e gabimit; përfundimin. Parametrat nuk duhen fshehur brenda funksioneve. Një funksion duhet të kryejë një detyrë të qartë dhe të kthejë rezultate që mund të testohen.

\begin{table}[ht]
\centering
\caption{Dy rrugë teknike me të njëjtin standard shkencor.}
\begin{tabularx}{\textwidth}{YY}
\toprule Python & MATLAB/Octave\\ \midrule
NumPy për vektorë dhe matrica & Vektorë dhe matrica si objekte bazë\\
Matplotlib për figura & \texttt{plot} dhe funksionet grafike\\
Funksione dhe module \texttt{.py} & Funksione dhe skripte \texttt{.m}\\
Jupyter Notebook & Jupyter me kernel Octave ose Live Script\\
\bottomrule
\end{tabularx}
\end{table}

\section{Verifikimi, vleftësimi dhe riprodhueshmëria}
Verifikimi pyet nëse kodi zgjidh saktë ekuacionet e zgjedhura. Vleftësimi pyet nëse ekuacionet përshkruajnë mjaftueshëm sistemin real për qëllimin e deklaruar. Verifikimi përdor zgjidhje analitike, ligje ruajtjeje, teste njësish dhe konvergjencë. Vleftësimi përdor matje, parametra të kalibruar dhe krahasim me modele alternative.

Riprodhueshmëria kërkon kodin, të dhënat, versionet e paketave, parametrat, farën e gjeneratorit dhe rendin e ekzekutimit. Një figurë pa këto të dhëna është dëshmi e paplotë.

\begin{lidhje}
Laboratori 1 ndërton të njëjtin notebook në njërën prej dy gjuhëve. Versioni i instruktorit shërben për shpjegim; versioni i studentit duhet të përmbajë arsyetimin, kodin dhe interpretimin e vet.
\end{lidhje}

\section*{Pyetje kontrolli}
\begin{enumerate}
\item Jepni një shembull ku kodi është i saktë, por modeli fizik është i papërshtatshëm.
\item Pse zvogëlimi i hapit nuk garanton gjithmonë përmirësim?
\item Cilat të dhëna duhen ruajtur për një eksperiment Monte Carlo të riprodhueshëm?
\end{enumerate}
\section{Algoritmi bazë i një eksperimenti numerik}
Algoritmi kryesor i kapitullit nuk është një formulë e vetme, por cikli i plotë i punës. Ai duhet të përdoret në çdo notebook të kursit.

\begin{enumerate}
\item Formulohet pyetja fizike dhe përcaktohet madhësia e daljes.
\item Shënohen supozimet, njësitë, parametrat dhe kushtet fillestare.
\item Zgjidhet modeli minimal që mund t'i përgjigjet pyetjes.
\item Përcaktohen rrjeti, algoritmi, toleranca dhe kriteri i ndalimit.
\item Shkruhet kodi si funksione të testueshme.
\item Verifikohet me një rast analitik ose një ligj ruajtjeje.
\item Kryhet studimi i konvergjencës ose i numrit të mostrave.
\item Raportohen rezultati, gabimi numerik, pacaktueshmëria dhe kufijtë e modelit.
\end{enumerate}

\subsection{Një skelet Python për notebook-un}
\begin{lstlisting}[language=Python,caption={Skelet minimal për një eksperiment numerik.}]
from dataclasses import dataclass
import numpy as np
import matplotlib.pyplot as plt

@dataclass(frozen=True)
class Parametra:
    y0: float = 1.5       # m
    v0: float = 8.0       # m/s
    g: float = 9.81       # m/s^2
    t_fund: float = 1.2   # s
    dt: float = 0.02      # s

def modeli_analitik(t, p):
    """Pozicioni vertikal pa rezistence ajri."""
    return p.y0 + p.v0*t - 0.5*p.g*t**2

def eksperiment(p):
    t = np.arange(0.0, p.t_fund + p.dt/2, p.dt)
    y = modeli_analitik(t, p)
    assert np.isclose(y[0], p.y0)
    return t, y

p = Parametra()
t, y = eksperiment(p)
plt.plot(t, y)
plt.xlabel("Koha t [s]")
plt.ylabel("Lartesia y [m]")
plt.grid(alpha=0.25)
plt.show()
\end{lstlisting}

\subsection{Kontrolli i riprodhueshmërisë}
Për simulime stokastike përdoret një gjenerator i deklaruar, jo gjendja globale e fshehur.
\begin{lstlisting}[language=Python]
rng = np.random.default_rng(2026)
mostra = rng.normal(loc=0.0, scale=1.0, size=10_000)
\end{lstlisting}
Fara nuk e bën rezultatin ``më të rastit''; ajo identifikon realizimin. Në raport shënohen fara, numri i mostrave dhe versionet e bibliotekave.

\subsection{Terminologji dhe stil shkencor}
Në këtë tekst përdoren \emph{zbatim} për implementation, \emph{rrjetë} për grid, \emph{hap} për step, \emph{mbetje} për residual, \emph{vleftësim} për validation dhe \emph{riprodhueshmëri} për reproducibility. Kodi nuk duhet të fshehë njësitë fizike; ato shënohen në komente, tabela dhe boshte.

\begin{lidhje}
Notebook-u i studentit duhet të nisë nga ky skelet dhe të ruajë ndarjen ndërmjet parametrave, modelit, eksperimentit dhe paraqitjes grafike.
\end{lidhje}

Për organizimin terminologjik të modelit, algoritmit dhe zbatimit numerik janë konsultuar tekstet universitare shqip të analizës numerike \cite{hanelli,hoxha,datja}, krahas literaturës ndërkombëtare të fizikës llogaritëse \cite{newman,landau}.

\subsection{Nga ligji fizik te algoritmi: një shembull i plotë}
Le të ndërtojmë modelin e rënies vertikale me rezistencë lineare të ajrit. Zgjedhim boshtin pozitiv për poshtë dhe shënojmë me $y(t)$ zhvendosjen dhe me $v(t)=\dot y(t)$ shpejtësinë. Ligji i dytë i Njutonit jep
\begin{equation}
 m\dot v=mg-cv,\qquad \dot y=v,
\end{equation}
ku $m>0$, $g>0$ dhe $c\ge0$. Përmasat fizike janë $[g]=LT^{-2}$ dhe $[c]=MT^{-1}$; prandaj $c/m$ ka përmasën e inversit të kohës. Koha karakteristike $\tau=m/c$ tregon sa shpejt humbet kujtesa e kushtit fillestar. Zgjidhja analitike,
\begin{equation}
 v(t)=v_\infty+(v_0-v_\infty)e^{-t/\tau},\qquad v_\infty=\frac{mg}{c},
\end{equation}
është një pikë reference për verifikimin e kodit. Ajo nuk përdoret sepse problemi është i vështirë, por sepse na lejon të ndajmë gabimin e algoritmit nga gabimi i modelit.

Me metodën e Eulerit, derivati në $t_n$ zëvendësohet nga diferenca përpara,
\begin{equation}
 \dot v(t_n)=\frac{v(t_n+h)-v(t_n)}{h}+\Oh(h).
\end{equation}
Duke neglizhuar termin $\Oh(h)$ dhe duke shënuar $v_n\approx v(t_n)$ merret
\begin{equation}
 v_{n+1}=v_n+h\left(g-\frac{c}{m}v_n\right),\qquad y_{n+1}=y_n+hv_n.
\end{equation}
Gabimi lokal i një hapi është $\Oh(h^2)$, ndërsa grumbullimi në një interval kohor të fiksuar jep gabim global $\Oh(h)$. Kjo dallon rendin lokal nga rendi global dhe shpjegon pse përgjysmimi i hapit duhet ta përgjysmojë afërsisht gabimin për një metodë të rendit të parë.

\begin{lstlisting}[language=Python,caption={Eksperiment konvergjence për rënien me rezistencë lineare.}]
import numpy as np

def renie_euler(m, c, g, y0, v0, T, h):
    n = int(np.ceil(T / h))
    h = T / n                 # hapi rregullohet që t_n të arrijë saktë T
    t = np.linspace(0.0, T, n + 1)
    y = np.empty(n + 1); v = np.empty(n + 1)
    y[0], v[0] = y0, v0
    for k in range(n):
        y[k + 1] = y[k] + h * v[k]
        v[k + 1] = v[k] + h * (g - (c / m) * v[k])
    return t, y, v

def v_sakte(t, m, c, g, v0):
    tau = m / c
    v_inf = m * g / c
    return v_inf + (v0 - v_inf) * np.exp(-t / tau)

for h in (0.2, 0.1, 0.05, 0.025):
    t, y, v = renie_euler(1.0, 0.4, 9.81, 0.0, 0.0, 4.0, h)
    gabimi = np.max(np.abs(v - v_sakte(t, 1.0, 0.4, 9.81, 0.0)))
    print(f"h={h:7.4f}, gabimi maksimal={gabimi:.6e}")
\end{lstlisting}
Ky program ndan modelin, integruesin dhe testin analitik. Në MATLAB/Octave të njëjtat vargje ndërtohen me \texttt{linspace} dhe cikli me \texttt{for}; ndryshimi sintaksor nuk ndryshon logjikën shkencore.

\subsection{Verifikim, vleftësim dhe analizë ndjeshmërie}
Një rezultat mund të jetë i konverguar ndaj $h\to0$ dhe përsëri të jetë fizikisht i gabuar. Konvergjenca kontrollon vetëm diskretizimin. Për vleftësim duhen të dhëna ose një model referimi. Në shembullin e mësipërm, koeficienti $c$ mund të varet nga forma, dendësia e ajrit dhe regjimi i rrjedhjes. Nëse $c$ është matur me pacaktueshmëri, kjo pacaktueshmëri duhet të përhapet deri te $v(t)$.

Ndjeshmëria lokale ndaj parametrit $c$ mund të vlerësohet me
\begin{equation}
 S_c(t)\approx\frac{v(t;c+\Delta c)-v(t;c-\Delta c)}{2\Delta c}.
\end{equation}
Vlera $S_c$ ka njësi dhe varet nga shkalla; një ndjeshmëri pa përmasa është $cS_c/v$. Në një studim të plotë raportohen veçmas gabimi i diskretizimit, pacaktueshmëria e parametrave dhe mospërputhja model--eksperiment.

\subsection{Përmbledhje dhe punë e pavarur}
Eksperimenti numerik është një zinxhir argumentesh: përkufizim i pyetjes, zgjedhje e modelit, kontroll dimensional, diskretizim, implementim, verifikim dhe interpretim fizik. Studenti duhet të jetë në gjendje të tregojë ku hyn çdo supozim dhe cilin lloj gabimi kontrollon çdo provë.

\begin{enumerate}
\item Nxirrni zgjidhjen analitike për $y(t)$ dhe përdoreni për të matur rendin e Eulerit.
\item Zëvendësoni forcën $-cv$ me $-cv|v|$ dhe përcaktoni shpejtësinë kufitare.
\item Ndërtoni të njëjtin eksperiment në MATLAB/Octave dhe krahasoni rezultatet deri në tolerancën e rrumbullakimit.
\item Shpjegoni pse krahasimi me të dhënat eksperimentale është vleftësim, ndërsa krahasimi me zgjidhjen analitike është verifikim.
\end{enumerate}

\subsection{Analiza pa përmasa dhe projektimi i simulimit}
Para se të zgjidhet një metodë numerike, është e dobishme të identifikohen shkallët karakteristike të problemit. Në modelin e rënies me rezistencë lineare,
\begin{equation}
 m\frac{\dd v}{\dd t}=mg-cv,
\end{equation}
shpejtësia karakteristike është $V=mg/c$ dhe koha karakteristike është $T=m/c$. Duke vendosur $u=v/V$ dhe $\tau=t/T$, ekuacioni bëhet
\begin{equation}
 \frac{\dd u}{\dd\tau}=1-u.
\end{equation}
Të tre parametrat dimensionalë janë reduktuar në një ekuacion universal. Kjo formë tregon se çdo trup që i bindet të njëjtit model ka të njëjtën kurbë $u(\tau)$; parametrat vetëm ndryshojnë shkallët e boshteve. Në simulim, analiza pa përmasa ndihmon në zgjedhjen e hapit: $\Delta\tau=\Delta t/T$ duhet të jetë mjaft i vogël për të përshkruar relaksimin. Ajo ndihmon edhe në testim, sepse rezultatet për grupe të ndryshme parametrash duhet të bien mbi të njëjtën kurbë të reduktuar.

Në modelin kuadratik $m\dot v=mg-cv|v|$, shpejtësia kufitare është $V=\sqrt{mg/c}$ dhe një shkallë kohe është $T=V/g$. Përsëri merret një ekuacion pa përmasa, por struktura jolineare mbetet. Krahasimi i dy modeleve duhet bërë në të njëjtat madhësi të reduktuara; përndryshe ndryshimi i njësive ose i shkallëve mund të ngatërrohet me ndryshim fizik.

\subsubsection{Nga kërkesa shkencore te kriteri i pranimit}
Një simulim duhet të fillojë me një kërkesë të matshme. Pohimi ``të simulohet lavjerrësi'' nuk përcakton as produktin, as saktësinë. Një kërkesë operative mund të jetë: ``të përcaktohet perioda për $0\le\theta_0\le1.2$ rad me gabim relativ nën $10^{-3}$ dhe të kontrollohet ruajtja e energjisë për njëqind perioda''. Kjo fjali përcakton madhësinë dalëse, domenin e parametrave, tolerancën dhe një invariant verifikimi.

Kriteri i pranimit nuk duhet të jetë vetëm ``kodi ekzekutohet''. Për një metodë të rendit $p$, raporti
\begin{equation}
 R_h=\frac{\|y_h-y_{h/2}\|}{\|y_{h/2}-y_{h/4}\|}
\end{equation}
duhet të afrohet me $2^p$ në regjimin asimptotik. Nëse kjo nuk ndodh, mund të ketë gabim implementimi, hap ende të madh, zgjidhje jo të lëmuar ose dominim të rrumbullakimit. Kriteret fizike përfshijnë shenjat, kufijtë, njësitë dhe ligjet e ruajtjes. Kriteret statistike përfshijnë përsëritjen me fara të ndryshme dhe intervalin e besimit.

\subsubsection{Hierarkia e pacaktueshmërive}
Në një raport të plotë dallohen së paku katër nivele. Pacaktueshmëria e hyrjeve vjen nga matjet e parametrave. Gabimi i modelit vjen nga mekanizmat e neglizhuar. Gabimi numerik vjen nga diskretizimi dhe aritmetika me saktësi të fundme. Gabimi i kampionimit vjen nga përdorimi i një numri të fundmë realizimesh. Zvogëlimi i hapit kontrollon vetëm një pjesë të nivelit të tretë; rritja e numrit të mostrave kontrollon vetëm nivelin e katërt. Një grafik me shumë shifra nuk mund të kompensojë një model fizik të papërshtatshëm.

Për laboratorin rekomandohet që çdo notebook të përmbajë një tabelë të shkurtër me kolonat: burimi i pacaktueshmërisë, mënyra e vlerësimit, madhësia e efektit dhe prova e kontrollit. Kjo e kthen notebook-un nga një demonstrim kodi në dokument shkencor.

